\section{Tensoring and Flatness}

\subsection{Tensor products}

I assume some familiarity with the tensor product (see e.g.\ my summary on representation theory otherwise).
The tensor product is a bifunctor
$$ \bul\otimes_B\bul\colon {}_A\Mod_B\times{}_B\Mod_C \to {}_A\Mod_C , $$
with the universal property that it represents the balanced map functor:
$$ \hom(M\otimes_BN,P) \cong {\rm Bal}(M,N;\,P) . $$
Balanced maps are maps $M\times N\to P$ which are additive and multiplicative in each coordinate, and $f(mb,n)=f(m,bn)$ for all $b\in B$.

The tensor product also has a right-adjoint.
For an $(A,B)$-bimodule $M$, $\bul\otimes_AM$ is a functor $\Mod_A\to\Mod_B$ with right adjoint $\hom_B(M,\bul)$.

In the case of commutative rings (this course), all modules are bimodules, so the tensor product is a bifunctor
$$ \bul\otimes_A\bul\colon \Mod_A\to\Mod_A, $$
which represents the bilinear map functor (bilinear in the usual sense).

Note still that we use tensor products between module categories.
In particular if $B$ is an $A$-algebra, then $B\otimes_AM$ is a $B$-module ($b(c\otimes m)=(bc)\otimes m$).
This defines a map $\Mod_A\to\Mod_B$, and its right-adjoint is the restriction map (restrict a $B$-module to an $A$-module).
Note that this is a special case of the tensor-hom adjunction: $\bul\otimes_AB$ is left-adjoint to $\hom_B(B,\bul)\cong B$.

So we have a natural isomorphism
$$ \hom_B(B\otimes_AM,P) \cong \hom_A(M,P|_A),\qquad M\in\Mod_A, P\in\Mod_B . $$
The correspondence is given by mapping $f\colon B\otimes_AM\to P$ to $\tilde f\colon M\to P$ by $\tilde f(m)=f(1\otimes m)$.
And $g\colon M\to P|_A$ to $\hat g\colon B\otimes_AM\to P$ by $\hat g(1\otimes m)=g(m)$.

\bprop

    Let $M$ be an $A$-module.
    \benum
        \item For a multiplicative set $S\subseteq A$, $S^{-1}A\otimes_AM$ is naturally isomorphic to $S^{-1}M$.
        \item For an ideal $I\leq A$, $A/I\otimes_AM$ is naturally isomorphic to $M/IM$.
    \eenum

\eprop

\bproof

    The canonical map $M\to S^{-1}M$ mapping $m\mapsto m/1$ defines a map $f\colon S^{-1}A\otimes_AM\to S^{-1}M$ by
    $f(1\otimes m)=m/1$.
    Then
    $$ f(a/s\otimes m) = a/sf(1\otimes m) = am/s . $$
    We define its inverse $g\colon S^{-1}M\to S^{-1}A\otimes_AM$ by
    $$ g(m/s) = 1/s\otimes m . $$
    We claim this is well-defined: if $m/s=m'/s'$ then $1/s\otimes m=1/s'\otimes m'$.
    If $t(ms'-m's)$ then
    $$ \frac1s\otimes m = \frac{ts'}{ts's}\otimes m = \frac1{ts's}\otimes ts'm = \frac1{ts's}\otimes tsm' = \frac1{s'}\otimes m' $$
    as desired, and clearly $g$ is a module morphism.
    $f$ and $g$ are also clearly morphisms.

    The proof for the second point is similar: projection $M\to M/IM$ is an $A$-module morphism.
    Thus this gives rise to an $A/IA$-module morphism $f\colon(A/IA)\otimes_AM\to M/IM$ by $f(1\otimes m)=\bar m$.
    Then
    $$ f(\bar a\otimes m) = af(1\otimes m) = a\bar m = \bar{am} . $$
    Its inverse is $g\colon M/IM\to(A/IA)\otimes_AM$ given by $g(\bar m)=1\otimes m$; this is easily seen to be well-defined and an inverse.
    \qed

\eproof

\bcoro

    Let $B$ be an $A$-algebra and $M$ a $B$-module.
    Then there is a natural map
    $$ B\otimes_A M|_A \to M,\qquad b\otimes m\to bm . $$
    In particular this is a natural map $m_B\colon B\otimes_AB\to B$ by $b\otimes b'\mapsto bb'$.

\ecoro

\bproof

    This map is the map corresponding to the identity of $\hom_A(M|_A,M|_A)$.
    The second point is for $M=B$.
    \qed

\eproof

\blemm

    If $B,C$ are $A$-algebras, then their tensor product $B\otimes_AC$ is naturally an $A$-algebra and their coproduct in the category of
    $A$-algebras.

\elemm

\bproof

    Let $D=B\otimes_AC$.
    The natural algebra structure is given by $(b\otimes c)(b'\otimes c')=(bb')\otimes(cc')$, and $A\to D$ by $a\mapsto a(1\otimes 1)$.
    This is the composition of
    $$ D\times D\to D\otimes_AD,\qquad (d,d')\mapsto d\otimes d', $$
    and
    $$ D\otimes_AD = (B\otimes_AC)\otimes_A(B\otimes_AC) \cong (B\otimes_AB)\otimes_A(C\otimes_AC) \to B\otimes_AC $$
    where the final map is the tensor of $m_B\otimes m_C$.

    The inclusion maps $i_B\colon B\to D$ and $i_C\colon C\to D$ are given by $b\mapsto b\otimes1$ and $c\mapsto1\otimes c$.

    Now if $E$ is another $A$-algebra and there are maps $f\colon B\to E,g\colon C\to E$, then if $h\colon D\to E$ commutes with the
    rest of the diagram, i.e.
    $$ g = h\circ i_C,\qquad f = h\circ i_B, $$
    then we must have
    $$ h(b\otimes c) = h((b\otimes1)(1\otimes c)) = h(b\otimes1)h(1\otimes c) = hi_B(b)hi_C(c) = f(b)g(c) . $$
    Thus $h$ is unique, and this is well-defined (it is bilinear).
    \qed

\eproof

\subsection{Flatness}

\bdefn

    A functor is {\emphcolor left-exact} if it preserves all finite limits, and is {\emphcolor right-exact} if it preserves all finite colimits.
    A functor is {\emphcolor exact} if it is both left- and right-exact (preserves all finite limits and colimits).

\edefn

Clearly the composition of (left/right-) exact functors are (left/right-) exact.

We quickly note a few equivalences which are useful to keep in mind.
The proof provided is not complete, but is not too complicated to complete.

\bprop

    For an additive functor $F$ between Abelian categories (in particular between module categories), the following are equivalent:
    \benum
        \item $F$ is left-exact,
        \item If $0\to A\to B\to C$ is an exact sequence, then the induced sequence $0\to FA\to FB\to FC$ is exact as well,
        \item If $0\to A\to B\to C\to0$ is a short exact sequence, then the induced sequence $0\to FA\to FB\to FC$ is exact.
    \eenum
    In particular from (2) we see that left-exact functors preserve injections.

    Similarly, the following are equivalent:
    \benum
        \item $F$ is right-exact,
        \item If $A\to B\to C\to0$ is an exact sequence, then the induced sequence $FA\to FB\to FC\to0$ is exact as well,
        \item If $0\to A\to B\to C\to0$ is a short exact sequence, then the induced sequence $FA\to FB\to FC\to0$ is exact.
    \eenum
    In particular right-exact functors preserve surjections.

    Consequently, the following are equivalent:
    \benum
        \item $F$ is exact,
        \item $F$ preserves short exact sequences,
        \item $F$ preserves long exact sequences,
        \item $F$ is left-exact and preserves surjections,
        \item $F$ is right-exact and preserves injections.
    \eenum

\eprop

\bproof

    First we note that an additive functor immediately preserves all finite biproducts.
    And $0\to A\to B\to C$ being exact is equivalent to it being a kernel diagram.
    Thus $F$ is left-exact iff it preserves it.
    And $F$ preserves left-exact sequences (i.e.\ ones in point (2)) iff it maps short exact sequences to left-exact ones.
    Clearly if it preserves left-exact sequences it maps short exact ones to left-exact ones (SES are in particular left-exact),
    and the converse is shown by restricting the rightmost map to its image.

    Since all injections are kernels, we see that left-exact functors preserve injections.

    Similarly we show the equivalences for right-exact functors.

    By the previous two equivalences, we see $F$ being exact is equivalent to preserving SES sequences.
    And preserving long-exact sequences is done by properly restricting and quotienting.
    If $F$ is left-exact and preserves surjections, then if $0\to A\to B\to C\to0$ is exact, we know that $0\to FA\to FB\to FC$ is exact,
    and $FB\to FC\to0$ is exact (as it preserves surjections).
    Putting this together, if $F$ is left-exact and preserves surjections it is exact.
    \qed

\eproof

\bprop

    Left-adjoint functors are right-exact, and right-adjoint functors are left-exact.

\eprop

\bproof

    Left-adjoint functors preserve {\it all\/} colimits, and are thus in particular right-exact.
    Similar for right-adjoint functors.
    \qed

\eproof

Tensor products are left-adjoints and thus right-exact:

\bprop

    If $M'\to M\to M''\to0$ is an exact sequence of $A$-modules, then for any $A$-module $N$, the induced sequence
    $$ M'\otimes_AN \to M\otimes_AN \to M''\otimes_AN \to 0 $$
    is exact.

\eprop

\bnote

    Tensor products are not exact; they do not preserve injections.
    Indeed, if $f\colon A\to A$ is the map $f(x)=ax$ for $a\in A$, then $f\otimes M=f\otimes\id_M$ is the endomorphism of
    $A\otimes_AM\cong M$ which maps $m\mapsto am$.
    When $a$ is not a zero divisor (e.g.\ $A$ is a domain), $f$ is injective, but $f\otimes M$ need not be (e.g.\ $M=A/(a)$).

\enote

Since $\hom_A(M,\bul)$ is right-adjoint, it is left exact so it preserves finite limits.
It is not hard to show that $\hom_A(\bul,M)$ is left-exact as well (but contravariantly).
Furthermore the Yoneda embedding reflects exactness:

\bprop

    A sequence of $A$-modules $M_\bul=M'\to M\to M''\to0$ is exact iff
    $$ \hom(M_\bul,N) = 0 \to \hom(M'',N) \to \hom(M,N) \to \hom(M',N) $$
    is exact for all $N\in\Mod_A$.
    Similarly $N_\bul=0\to N'\to N\to N''$ is exact iff
    $$ \hom(M,N_\bul) = 0 \to \hom(M,N') \to \hom(M,N) \to \hom(M,N'') $$
    is exact for all $M$.

\eprop

\bproof

    The ``only if'' part is due to left exactness.
    The converse (second point) is done by choosing $M=N'$ and looking at where $\id$ goes (this shows that the composition of maps
    is zero), then $M$ to be the kernel of the second map, and showing that this implies exactness at $M$, and then you just need to
    show the first map is injective.
    \qed

\eproof

While not all tensor products are exact, some are.
We give this case a special term.

\bdefn

    An $A$-module $M$ is {\emphcolor flat} if the functor $\bul\otimes_AM$ is exact.

\edefn

\bnote

    Since $\bul\otimes_AM$ is right-exact for any module $M$, $M$ is flat iff $\bul\otimes_AM$ preserves injections.

\enote

\bexam

    If $S\subseteq A$ is multiplicative, then $\bul\otimes_AS^{-1}A$ is naturally isomorphic to $S^{-1}\bul$, which is exact.
    Thus $S^{-1}A$ is a flat $A$-module.

    If $M$ is free, so $M\cong A^{\oplus I}$, then
    $$ \bul\otimes_AM \cong \bul\otimes_AA^{\oplus I} (\bul\otimes_AA)^{\oplus I} \cong \bul^{\oplus I} . $$
    The second isomorphism is due to left adjoints preserving colimits.
    And $\bul^{\oplus I}$ is exact (it maps injective maps to injective maps, more generally direct sums in module categories are exact,
    satisfying Grothendieck's axiom (AB4)).

\eexam

\blemm

    Let $M$ be an $A$-module and $S\subseteq A$ a multiplicative set.
    \benum
        \item If $N$ is an $S^{-1}A$-module, it is an $A$-module via the canonical map $A\to S^{-1}A$, and there is a natural isomorphism
        $$ M\otimes_AN \to S^{-1}M\otimes_{S^{-1}A}N,\qquad m\otimes n\mapsto m/1\otimes n . $$
        \item For an $A$-module $N$, there is a natural isomorphism
        $$ S^{-1}M\otimes_{S^{-1}A}S^{-1}N \to S^{-1}(M\otimes_AN),\qquad \frac ms\otimes\frac nt \mapsto \frac{m\otimes n}{st} . $$
        In particular for every prime $\p\leq A$, there is a natural isomorphism
        $$ M_\p \otimes_{A_\p} N_\p \cong (M\otimes_AN)_\p . $$
    \eenum

\elemm

\bproof

    There is a natural isomorphism (by associativity of the tensor)
    $$ S^{-1}M\otimes_{S^{-1}A}N \cong (M\otimes_AS^{-1}A)\otimes_{S^{-1}A}N \cong M\otimes_AN $$
    which maps $m/s\otimes n$ to $m\otimes 1/s\otimes n$ to $m\otimes n/s$.
    So in reverse it maps $m\otimes n$ to $m/1\otimes n$ as desired.

    There is a natural isomorphism
    $$ S^{-1}M\otimes_{S^{-1}A}S^{-1}N \cong M\otimes_AS^{-1}N \cong M\otimes_A(N\otimes_AS^{-1}A) \cong (M\otimes_AN)\otimes_A S^{-1}A , $$
    and this is naturally isomorphic to $S^{-1}(M\otimes_AN)$.
    \qed

\eproof

\bcoro

    Flatnes is local, i.e.\ the following are equivalent:
    \benum
        \item $M$ is a flat $A$-module,
        \item $S^{-1}M$ is a flat $S^{-1}A$-module for all multiplicative $S\subseteq A$,
        \item $M_\p$ is a flat $A_\p$-module for all prime $\p\leq A$,
        \item $M_\m$ is a flat $A_\m$-module for all maximal $\m\leq A$.
    \eenum

\ecoro

\bproof

    If $M$ is flat then by the above lemma
    $$ \bul\otimes_{S^{-1}A}S^{-1}M \cong S^{-1}(\bul\otimes_AM) , $$
    which is a composition of exact functors and thus exact.

    Now if $M_\m$ is flat for all maximal $\m\leq A$, we show $M$ is flat.
    Recall we need only show that $\bul\otimes_AM$ preserves injectives.
    Recall that $f$ is injective iff $f_\m$ is injective for all maximal $\m\leq A$.
    Now, $f_\m$ is isomorphic to $f\otimes M_\m$, which is injective since $M_\m$ is flat, we see $f_\m$ is injective for all $\m$ as desired.
    \qed

\eproof

\subsection{The Tor functor}

We now quickly state a few definitions and results from homological algebra.

\bdefn

    Let $\A,\B$ be Abelian categories (e.g.\ module categories).
    Then a {\emphcolor (homological) $\delta$-functor} is a collection of functors $\set{F_n\colon\A\to\B}_{n\geq0}$ satisfying the following two properties:
    \benum
        \item If $0\to A\to B\to C\to0$ is an SES in $\A$, then there exist maps $\delta_n\colon F_nC\to F_{n-1}A$ inducing a long exact sequence:
        $$ \cdots\to F_{n+1}C \xtos\delta F_nA \to F_nB \to F_nC \xtos\delta F_{n-1}A \to \cdots . $$
        \item The maps $\delta$ are natural: if we have a map of SES $0\to A\to B\to C\to0$ to $0\to A'\to B'\to C'\to0$, the following diagram commutes:

        \commsq{F_nC&F_{n-1}A\cr F_nC'&F_{n-1}A'}
            {\delta}{\delta}{}{}
    \eenum
    The functor $F_0$ is called the {\emphcolor base} of the $\delta$-functor.
    Cohomological $\delta$-functors are defined similarly but with the arrows inverted.

    If $F$ and $G$ are $\delta$-functors, then a morphism of $\delta$-functors is a sequence of natural transformations $f_n\colon F_n\To G_n$ which commute
    with the $\delta$-maps.

    A {\emphcolor universal} $\delta$-functor is a $\delta$-functor $F$ such that for any other $\delta$-functor $G$ and a natural transformation $f_0\colon G_0\to T_0$,
    $f_0$ extends uniquely to a morphism of $\delta$-functors $G\to T$.

\edefn

Clearly universal $\delta$-functors are unique up to isomorphism.
A functor $F$ is called {\emphcolor coeffaceable} if for every object $A$ there is a surjective $g\colon B\to A$ such that $F(g)=0$.
A $\delta$-functor $F$ is {\emphcolor coeffaceable} if every $F_n$ is coeffaceable for $n\geq1$.
An important due to Grothendieck is that a coeffaceable $\delta$-functor is universal.

Notice that if $F$ is a $\delta$-functor then its base $F_0$ is right-exact.
This is because if $0\to A\to B\to C\to0$ is exact, then the end of the long exact sequence induced by $F$ is
$$ \cdots\to F_0A \to F_0B \to F_0C \to 0 . $$
Another result, due to Grothendieck again, is that every right-exact functor is the base of a (necessarily unique up to isomorphism) universal $\delta$-functor (under some
basic assumptions on the underlying category).
This is called the {\emphcolor derived functor}.

In particular, $\bul\otimes_AM$ is right-exact for any $A$-module $M$.
Thus it has a derived functor, called the {\emphcolor Tor-functor}, denoted by $\tor_n^A(\bul,M)$.
We can also compute the derived functor of $N\otimes_A\bul$ giving us another functor $\tor_n^A(N,\bul)$, and a non-trivial result tells us that these are two
parts of the same bifunctor $\tor_n^A(\bul,\bul)\colon\Mod_A\times\Mod_A\to\Mod_A$.

Since the tensor product is commutative (up to isomorphism), the derived functors $\tor_n^A(M,\bul)$ and $\tor_n^A(\bul,M)$ are isomorphic (by uniqueness of
universal $\delta$-functors).

We summarize these results in the following theorem:

\bthrm[title=The Tor functor]

    There is a unique (up to unique isomorphism) collection of additive bi-functors $\tor_n^A\colon\Mod_A\times\Mod_A\to\Mod_A$ such that
    \benum
        \item $\tor_0^A=\bul\otimes_A\bul$,
        \item $\tor_n^A(\bul,M)$ is a $\delta$-functor for each $A$-module $M$.
        \item $\tor_n^A(N,M)\cong\tor_n^A(M,N)$ for all $A$-modules $N,M$ and all $n\geq0$.
    \eenum

\ethrm

The construction of derived functors is relatively straight-forward, but proving that they are (1) functors, (2) $\delta$-functors, (3) universal $\delta$-functors
is less trivial.
(1) and (2) is done pretty much directly as follows.
Suppose $F$ is right-exact, and let $A$ be an object (a module in our case).
Then we define a {\emphcolor projective resolution} of $A$ to be a long exact sequence
$$ P_\bul\to A = \cdots \xtos{d_1} P_1 \xtos{d_0} P_0 \to A \to 0 , $$
where each $P_n$ is a {\emphcolor projective object}.
In category of modules, free modules are a type of projective object (but not in general the only type), so we can instead consider only resolutions by free modules.

We then consider the induced chain complex
$$ F(P_\bul) = \cdots \xtos{Fd_2} F(P_2) \xtos{Fd_1} F(P_1) \xtos{Fd_0} F(P_0) \to 0 . $$
And then we define $F_n$ to be the $n$th homology modules of this sequence:
$$ F_n(A) = H_n(F(P_\bul)) = \frac{\ker Fd_n}{\im Fd_{n+1}} . $$
It can be shown that any map $f\colon A\to B$ there is a unique lift of $f$ to a sequence of maps between projective resolutions of $A$ and $B$ (unique up to homotopy),
and we define $F_n(f)$ to be the induced map on the homology modules.

It can be shown that this construction is not dependent on the choice of projective resolution: all projective resolutions give naturally isomorphic functors.
Furthermore, this is a $\delta$-functor.
And it is not hard to see that $F_0\cong F$.

Now notice that if $P$ is a projective object (whatever that means), certainly $0\to P\to P\to0$ is a projective resolution (the non-trivial map is the identity).
Thus $F_\bul$ are the homology modules of the chain complex $0\to FP\to0$, which vanish for $n\geq1$.
In particular, as we said that free modules are projective, if $N$ is a free $A$-module then $\tor_n^A(N,M)=0$ for all $M\in\Mod_A$.

Now we call an object $P$ {\emphcolor $F$-acylic} if $F_nP=0$ for all $n\geq1$, so clearly projectives are $F$-acylic for any $F$.
A simple result is that derived functors can be computed using $F$-acylic resolutions.
That is, if $P_\bul\to A\to0$ is a resolution of $A$ using $F$-acylic objects, then $F_n(A)\cong H_n(F(P))$.
Of course to know what objects are $F$-acylic we need to compute $F_n$ via usual projective resolutions.

But in our case, if $N\in\Mod_A$ is flat, then it is $\bul\otimes_AM$-acylic for all $M\in\Mod_A$ and vice versa.

\bprop

    Let $N\in\Mod_A$, then the following are equivalent:
    \benum
        \item $N$ is flat,
        \item $\tor^A_n(N,M)=0$ for all $M\in\Mod_A$ and all $n\geq1$,
        \item $\tor^A_1(N,M)=0$ for all $M\in\Mod_A$
    \eenum
    In particular free modules are flat.

\eprop

\bproof

    Consider a projective resolution of $M$, $P_\bul\to M\to0$.
    Then $\tor^A_n(N,M)$ are the homology modules of the chain complex
    $$ \cdots\to P_1\otimes_AN \to P_0\otimes_AN \to 0 . $$
    But $N$ is flat, and thus tensoring with it is exact so it preserves exact sequences.
    So the above chain complex is exact at every index $n\geq1$ (at $n=0$ we have taken away $M\otimes_AN$ so it is no longer exact, this is because we want $\tor^A_0\cong\bul\otimes_A\bul$).
    Thus the homology modules are all trivial, so $\tor^A_n(N,M)=0$ for all $n\geq1$; so we have shown that (1) implies (2).

    And (2) clearly implies (3).

    Now suppose (3) then let $0\to X\to Y\to Z\to0$ be an SES.
    Then by $\delta$-functorality there is a long exact sequence
    $$ \cdots \to \tor^A_1(N,Z) \to N\otimes_AX \to N\otimes_AY \to N\otimes_AZ \to 0 . $$
    Since $\tor^A_1(N,Z)=0$ by assumption, this shows that
    $$ 0 \to N\otimes_AX \to N\otimes_AY \to N\otimes_AZ \to 0 $$
    is exact, and thus $N$ is flat.

    And as explained, $\tor^A_n(P,\bul)=0$ for any projective, and thus free, module $P$.
    So by the equivalence, free modules are flat.
    \qed

\eproof

Notice that free modules are flat for a simpler reason, which we show in the following lemma.

\blemm

    The direct sums and filtered colimits of flat modules are flat.

\elemm

\bproof

    Since the tensor product is a left-adjoint it commutes with (preserves) colimits.
    That is,
    $$ \bul\otimes_A\bigoplus_iM_i \cong \bigoplus_i\bul\otimes_AM_i,\qquad \bul\otimes_A\coflim M_i \cong \coflim(\bul\otimes_AM_i) . $$
    Since direct sums and filtered colimits are exact (Grothendieck's axioms (AB4) and (AB5)), if each $M_i$ is flat and thus $\bul\otimes_AM_i$ are all exact,
    then the above are exact as well.
    \qed

\eproof

Since $A$ is a flat $A$-module (since $\bul\otimes_AA$ is isomorphic to the identity functor), direct sums of $A$ are flat.
Direct sums of $A$ are just free modules.

\blemm

    Homology commutes with exact functors: if $C_\bul$ is a chain complex and $F$ an exact functor, then
    $$ H_n(F(C_\bul)) \cong F(H_n(C_\bul)) \hbox{ for all $n$} . $$

\elemm

\bproof

    Notice that a quotient $M/N$ is the cokernel of the inclusion $N\inj M$, and the image of $f$ is just $\im f=\ker\coker f$.
    Since $F$ is exact, it preserves injections and cokernels, so $F(N)$ remains a subobject of $F(M)$ and $F(M)/F(N)\cong F(M/N)$ since $F$ preserves cokernels.
    And it preserves images as well since $F\im f=F\ker\coker f\cong\ker\coker Ff=\im Ff$.

    In particular,
    $$ F(H_n(C_\bul)) = F(\ker f_n/\im f_{n+1}) \cong (F\ker f_n)/(F\im f_{n+1}) \cong (\ker Ff_n)/(\im Ff_{n+1}) = H_n(F(C_\bul)) $$
    as desired.
    \qed

\eproof

\bthrm

    The Tor functors commute with direct sums and filtered colimits:
    $$ \tor^A_n\parens{N,\bigoplus_iM_i} \cong \bigoplus_i\tor^A_n(N,M_i),\qquad \tor^A_n(N,\coflim M_i) \cong \coflim\tor^A_n(N,M_i) . $$

\ethrm

\bproof

    We explained that Tor can be computed via acylic resolutions, and that flat modules are acylic.
    So if
    $$ P^i_\bul \to M_i \to 0 $$
    are flat resolutions of $M_i$ for each $i$, then by exactness of direct sums and filtered colimits,
    $$ \bigoplus_i P^i_\bul \to \bigoplus_iM_i \to 0,\qquad \coflim P^i_\bul \to \coflim M_i \to 0 $$
    are exact sequences.
    Furthermore we just showed that these are flat, and thus Tor can be computed using their homology.
    And again since direct sums and filtered colimits are exact, homology commutes with them.
    Putting everything together,
    $$ \eqalign{
        \tor^A_n\parens{N,\bigoplus_iM_i} &= H_n\parens{\bigoplus_iP^i_\bul} = \bigoplus_iH_n(P^i_\bul) = \bigoplus_i\tor^A_n(M_i),\cr
        \tor^A_n\parens{N,\coflim M_i} &= H_n\parens{\coflim P^i_\bul} = \coflim H_n(P^i_\bul) = \coflim \tor^A_n(M_i).
    } $$
    \qed

\eproof

Now that we have the basic properties of the Tor functor down, we will use it to prove properties about flatness.

\blemm

    An $A$-module $M$ is flat iff $\tor^A_1(M,A/I)=0$ for all ideals $I\leq A$.

\elemm

\bproof

    The ``only if'' case is clear, as the first torsion module of a flat module with any other module is zero.
    Conversely, we need to show that $\tor^A_1(M,N)=0$ for any $N\in\Mod_A$.

    First suppose that $N$ is fg, and we induct on the number of generators $a_1,\dots,a_n$.
    For $n=0$ it is clear since then $N=0$.
    And for $n=1$ we have that $N=(a_1)\cong A/\ann(a_1)$ and thus the result follows from the assumption.
    Now inductively consider the SES
    $$ 0 \to (a_n) \to N \to N/(a_n) \to 0 . $$
    Thus we have an exact sequence
    $$ \cdots \to \tor^A_1(M,(a_n)) \to \tor^A_1(M,N) \to \tor^A_1(M,N/(a_n)) \to \cdots . $$
    Now $N/(a_n)$ is generated by $n-1$ elements, and thus $\tor^A_1(M,N/(a_n))=0$ by induction, and by our base case $\tor^A_1(M,(a_n))=0$.
    Thus $\tor^A_1(M,N)=0$ by exactness.

    And if $N$ is not fg, then $N=\coflim N_i$ where $N_i$ ranges over all the fg submodules of $N$.
    Since $\tor$ commutes with filtered colimits, we have
    $$ \tor^A_1(M,\coflim N_i) \cong \coflim\tor^A_1(M,N_i) = \coflim 0 = 0 , $$
    as desired.
    \qed

\eproof

\bthrm

    An $A$-module $M$ is flat iff for every ideal $I\leq A$, the map $I\otimes_AM\to IM$ given by $a\otimes m\mapsto am$ is an isomorphism.

\ethrm

\bproof

    First note that the given map is induced by an $A$-bilinear map and is thus well-defined.

    If $M$ is flat and $I\leq A$ is an ideal, then we have a short exact sequence
    $$ 0 \to I \to A \to A/I \to 0 , $$
    where the first map is given by inclusion and the second by projection.
    By flatness, this induces an SES, recalling that $A\otimes_AM\cong M$,
    $$ 0 \to I\otimes_AM \to M \to M\otimes_A(A/I) \to 0 . $$
    The first map is then the map $M\otimes_AI\to M$ given by $m\otimes a\mapsto am$, the map we are investigating, which is injective by exactness.

    Recall that $M\otimes_A(A/I)\cong M/IM$ by $m\otimes(a+I)\mapsto am+IM$.
    Thus the second map corresponds to the map $M\to M/IM$ given by projection.
    By exactness, the image of the first map is the kernel of the second, which is $IM$.
    So we have that our map is injective onto its image, $IM$, and thus is an isomorphism $I\otimes_AM\to IM$ as desired.

    Conversely by the above lemma, it is sufficient to show that $\tor^A_1(M,R/I)=0$.
    From the SES $0\to I\to A\to A/I\to0$, we have an exact sequence
    $$ \tor^A_1(M,I) \to \tor^A_1(M,A) \to \tor^A_1(M,A/I) \to M\otimes_A I \to M \to M\otimes_A(A/I) \to 0 . $$
    Since $A$ is free, $\tor^A_1(M,A)=0$.
    And the map $M\otimes_AI\to M$ is precisely the map $m\otimes a\mapsto am$, which is injective by assumption.
    So we have an exact sequence
    $$ 0 \to \tor^A_1(M,A/I) \to M\otimes_AI \xtos f M , $$
    where $f$ is injective.
    This requires that $\tor^A_1(M,A/I)=0$ as desired.
    \qed

\eproof

\bexam

    Let $R=\bC[x]$ and $M$ be the complex algebra of functions $\bC\to\bC$.
    Then we define $R\to M$ by mapping $x$ to the zero map.
    This turns $M$ into an $R$-algebra, where $p\in R$ is mapped to $p(0)$.

    We claim that $M$ is not a flat $R$-module.
    Indeed, let $I=(x)$ so that $IM=0$.
    But $I\cong R$ as $R$-modules (since they are both free $R$-modules on a single generator), so $I\otimes_RM\cong M\neq0$.
    Thus $IM$ and $I\otimes_RM$ are not isomorphic and in particular $M$ is not flat.

\eexam

While free modules are always flat, the converse is not always true.
But it is in an important special case, whose proof is homework.

\bprop

    Suppose $A$ is a local ring and $M$ a fg $A$-module.
    Then $M$ is free iff it is flat.

\eprop

\bcoro

    Let $M$ be an $A$-module, then the following are equivalent:
    \benum
        \item $M$ is flat,
        \item $M_\p$ is a free $A_\p$-module for all prime $\p\leq M$,
        \item $M_\m$ is a free $A_\m$-module for all maximal $\m\leq M$.
    \eenum

\ecoro

\bproof

    Flatness is local: $M$ is flat iff $M_\p$ is a flat $A_\p$-module for all prime $\p\leq M$ or equivalently all maximal $\m\leq M$.
    And $A_\p$ is a local ring, so by the previous proposition this is iff $M_\p$ is free.
    \qed

\eproof

We can also give a nicer criteria in the case that our ring is a domain.

\blemm

    Let $A$ be a local domain with maximal ideal $\m$, and let $k=A/\m$ be its residue field and $K=\Frac(A)$ its field of fractions.
    Let $M$ be a fg $A$-module, and suppose $\dim_k(k\otimes_AM)=\dim_K(K\otimes_AM)=n$.
    Then $M$ is free of rank $n$.

\elemm

\bproof

    Nakayama's lemma implies that there is a generating set $x_1,\dots,x_n\in M$ such that their projections form a basis in $k\otimes_AM\cong M/\m M$
    (the final formulation of Nakayama's lemma).
    Define $\phi\colon A^n\to M$ which maps the standard basis vectors $e_i$ to $x_i$.
    This gives an SES
    $$ 0 \to \ker\phi \to A^n \to M \to 0 . $$
    Recall that localization of flat modules are flat, and since $K$ is the localization of $A$ by $A-0$, it is a flat $A$-module.
    Thus tensoring with $K$ preserves the SES,
    $$ 0 \to K\otimes_A\ker\phi \to K\otimes_AA^n \to K\otimes_AM \to 0 . $$
    Recall that tensoring commutes with direct sums, so this is isomorphic to
    $$ 0 \to K\otimes_A\ker\phi \to K^n \to K\otimes_AM \to 0 . $$
    By assumption $K\otimes_AM$ has dimension $n$ over $K$, and since the final map is a surjection it must be an isomorphism by considering dimensions.
    Thus $K\otimes_A\ker\phi=0$, and since localization commutes with tensoring,
    $$ 0 = K\otimes_A\ker\phi \cong \ker\phi_{A-0} . $$
    So for every $x\in\ker\phi$ we must have that $x/1=0$, i.e.\ there is an $a\in A-0$ such that $ax=0$, so $\ker\phi$ is a torsion module.
    But $\ker\phi\subseteq A^n$ which is not torsion as $A$ is a domain, so $\ker\phi=0$ and thus $\phi$ is an isomorphism and $M\cong A^n$ as desired.
    \qed

\eproof

\bthrm[title=Fiber dimension criterion for flatness]

    Let $A$ be a domain, and $M$ a fg $A$-module.
    Let $S=\spec(A)$, then $M$ is flat iff the map $S\to\bN$ given by $\p\mapsto\dim_{k(\p)}(k(\p)\otimes_AM)$ is constant.

\ethrm

\bproof

    As we showed by locality of flatness, $M$ is flat iff $M_\p$ is a free $A_\p$-module for all prime $\p\leq A$.
    So first suppose $M_\p$ is free for all prime $\p\leq A$, say $M_\p\cong A_\p^{n_\p}$ ($n_\p$ is finite since $M$ is fg).
    Then
    $$ k(\p) \otimes_A M = (A_\p/\p A_\p) \otimes_A M \cong M_\p/\p M_\p \cong A_\p^{n_\p}/(\p A_\p^{n_\p}) \cong k(\p)^{n_\p} , $$
    so $\dim_{k(\p)}(k(\p)\otimes_AM)=n_\p$.
    
    Now in general notice that if $\p\leq\q\leq A$ are prime, and $M$ is an $A$-module, then $M_\q$ remains an $A$-module and we can localize it by $\p$.
    It is easy to see that
    $$ (M_\q)_\p \cong M_\p . $$
    Thus,
    $$ M_{(0)} \cong (M_\p)_(0) \cong (A_\p^{n_\p})_{(0)} \cong (A_\p)_{(0)}^{n_\p} \cong A_{(0)}^{n_\p} = \Frac(A)^{n_\p} . $$
    And so we have that
    $$ \Frac(A)^{n_\p} \cong M_{(0)} \cong \Frac(A)\otimes_AM . $$
    The right-hand side is independent of $\p$, and thus $n_\p$ must be constant.
    Thus we see that $\p\mapsto n_\p=\dim_{k(\p)}(k(\p)\otimes_AM)$ is constant, as desired.

    Conversely suppose the map is constant, let its value be $n$.
    In particular we have that, for $K=\Frac(A)=A_{(0)}$,
    $$ n = \dim_{k(\p)}(k(\p)\otimes_AM) = \dim_K(K\otimes_AM) . $$
    We also have that
    $$ k(\p)\otimes_AM \cong (k(\p)\otimes_{A_\p}A_\p)\otimes_AM \cong k(\p) \otimes_{A_\p} (A_\p\otimes_AM) \cong k(\p) \otimes_{A_\p} M_\p . $$
    Thus we have
    $$ \dim_{k(\p)}(k(\p)\otimes_{A_\p}M_\p) = \dim_{k(\p)}(k(\p)\otimes_AM) = \dim_K(K\otimes_AM) . $$

    Now also we know
    $$ K\otimes_{A_\p}M_\p = (A_\p)_{(0)} \otimes_{A_\p} M_\p \cong (A_\p\otimes_{A_\p}M_\p)_{(0)} \cong (M_\p)_{(0)} \cong M_{(0)} \cong K\otimes_AM . $$
    Thus $\dim_K(K\otimes_AM)=\dim_K(K\otimes_{A_\p}M_\p)$, and so
    $$ \dim_{k(\p)}(k(\p)\otimes_{A_\p}M_\p) = \dim_K(K\otimes_{A_\p}M_\p) . $$
    By the previous lemma, $M_\p$ is free.
    \qed

\eproof

